% Generated draft: requires mathematical and source review.
\documentclass[11pt]{article}
\usepackage[margin=25mm]{geometry}
\usepackage{fontspec}
\setmainfont{Latin Modern Roman}
\usepackage{amsmath,amssymb,mathtools,unicode-math}
\setmathfont{Latin Modern Math}
\usepackage{enumitem,needspace,xurl,hyperref}
\setlength{\parindent}{0pt}
\setlength{\parskip}{5pt}
\setlength{\emergencystretch}{3em}
\newcommand{\N}{\mathbb{N}}
\newcommand{\Z}{\mathbb{Z}}
\newcommand{\Q}{\mathbb{Q}}
\newcommand{\R}{\mathbb{R}}
\newcommand{\C}{\mathbb{C}}
\newcommand{\F}{\mathbb{F}}
\newcommand{\A}{\mathbb{A}}
\newcommand{\PP}{\mathbb{P}}
\newcommand{\E}{\mathbb{E}}
\newcommand{\eps}{\varepsilon}
\newcommand{\dd}{\,\mathrm{d}}
\begin{document}
\section*{UnsolvedMath ID: 800390}
\textbf{Source problem code:} catalogue entry
\subsection{Definitions and mathematical statement}
Let \(X\) be a projective variety over \(\C\). An endomorphism \(f\colon X\to X\) is polarizable if there are an ample line bundle \(L\) on \(X\) and an integer \(q>1\) such that \(f^{*}L\cong L^{\otimes q}\). Write \(f^0=\operatorname{id}_X\). A point \(x\in X(\C)\) is preperiodic if \(f^m(x)=f^n(x)\) for some integers \(0\leq m<n\), and put
\[
\operatorname{PrePer}_f(X)=\{x\in X(\C):\exists\,0\leq m<n,\ f^m(x)=f^n(x)\}.
\]
A subvariety \(Y\subseteq X\) is preperiodic under \(f\) if \(f^m(Y)=f^n(Y)\) for some integers \(0\leq m<n\). Let \(\operatorname{Sing}(X)\) denote the singular locus, let \(T_{X,x}\) denote the Zariski tangent space at \(x\), and let \(\operatorname{Gr}_d(V)\) denote the Grassmannian of \(d\)-dimensional linear subspaces of a vector space \(V\).

\textbf{Conjecture (Ghioca--Tucker--Zhang, 2011; supplied as Conjecture 13 in the 2018 source list).} Suppose that \(f\colon X\to X\) is polarizable and that \(Y\subseteq X\) is a subvariety no irreducible component of which is contained entirely in \(\operatorname{Sing}(X)\). Set \(d=\dim Y\). The supplied source asserts the equivalence of the following conditions:
\begin{enumerate}
\item[(a)] \(Y\) is preperiodic under \(f\).
\item[(b)] There exists a subset
\[
S\subseteq Y(\C)\cap\operatorname{PrePer}_f(X),\qquad \overline{S}^{\mathrm{Zar}}=Y,
\]
consisting of smooth points such that, for every \(x\in S\), the tangent subspace \(T_{Y,x}\) is preperiodic under the source's stated induced action of \(f\) on
\[
\operatorname{Gr}_{d}(T_{X,x}).
\]
\end{enumerate}

\textbf{Unresolved qualifications in condition (b).} The supplied record does not fully specify the tangent-space action. The differential has the form
\[
\dd f_x\colon T_{X,x}\longrightarrow T_{X,f(x)},
\]
so it generally maps between different fibers, rather than acting on the single Grassmannian \(\operatorname{Gr}_d(T_{X,x})\). The natural candidate on pairs is
\[
(x,W)\longmapsto\bigl(f(x),\dd f_x(W)\bigr),
\]
where this expression belongs to the relevant Grassmannian only if \(\dim\dd f_x(W)=d\). For an everywhere-defined orbit of such pairs, preperiodicity would mean equality of two iterates of the pair for some \(0\leq m<n\). The supplied record does not specify a domain, rank conditions, or a convention at undefined iterates; this candidate therefore explains the missing data without adding it as a hypothesis.

The phrase ``smooth points'' also does not explicitly identify whether smoothness is required on \(Y\), on \(X\), or on both. Moreover, the record allows components of \(Y\) but does not explicitly require pure dimension. At a smooth point on a component of dimension less than \(d\), \(T_{Y,x}\) is not a point of \(\operatorname{Gr}_d(T_{X,x})\). Consequently, the displayed equivalence preserves the supplied conjectural formulation, but its tangent-action and dimension conventions require clarification before it is a fully specified mathematical assertion.
\subsection{Short English statement}
The supplied conjecture says that a subvariety of a complex projective variety with a polarizable endomorphism is preperiodic exactly when it contains a Zariski-dense set of preperiodic smooth points whose tangent subspaces are also preperiodic. No component of the subvariety may lie entirely in the ambient singular locus. The tangent action, smoothness requirement, and treatment of components of different dimensions remain unresolved in the supplied formulation.
\subsection{Sources}
\begin{itemize}
\item Local UnsolvedMath registry, supplied complete record with permanent ID \(800390\), exact source problem code \texttt{catalogue entry}, titled ``Open Problems in Arithmetic Dynamics --- Conjecture 13.'' The registry labels the record \texttt{open}; this is a supplied status label, not an independently verified mathematical or literature assessment.
\item Benedetto, DeMarco, Ingram, Jones et al., \emph{Current Trends and Open Problems in Arithmetic Dynamics} (2018), Conjecture 13. Supplied source URL: \url{https://arxiv.org/abs/1806.04980}. The record reports extraction by \texttt{arxiv\_latex} and an access date of 2026-10-16. These are supplied provenance details; no source visit or independent verification is asserted here.
\item The conjecture itself is attributed in the supplied statement to Ghioca--Tucker--Zhang (2011), reference \texttt{MR2854724}, Conjecture 2.4. The supplied record gives no further bibliographic details for this reference.
\end{itemize}
\end{document}
